% Part: first-order-logic
% Chapter: tableaux
% Section: rules-and-proofs

\documentclass[../../../include/open-logic-section]{subfiles}

\begin{document}

\iftag{FOL}
      {\olfileid{fol}{tab}{rul}}
      {\olfileid{pl}{tab}{rul}}

\olsection{القواعد و\usetoken{P}{tableau}}

ال!!^a{tableau} طريق منظّم لاستقصاء الوجوه التي قد تصدق بها
!!a{sentence} أو تكذب في \iftag{FOL}{!!a{structure}}{!!a{valuation}}.
ومادّة الجدول الدلالي ال!!{signed formula}s؛ وهي !!{sentence}s
تقرن كل واحدة منها بإشارة لقيمة الصدق، إما~$\True$ وإما~$\False$.
وتنتظم هذه ال!!{formula}s الموقّعة في شجرة تمتد من أعلاها إلى أسفلها.

\begin{defn}
  \emph{!!{signed formula}} هي زوج: أوله قيمة صدق، وثانيه
  !!a{sentence}. ولها إذن إحدى هاتين الصورتين:
  \[
  \sFmla{\True}{!A} \text{ أو } \sFmla{\False}{!A}.
  \]
\end{defn}

ولتصوّر المراد نقرأ~$\sFmla{\True}{!A}$: ``يجوز أن تكون~$!A$
صادقة''، ونقرأ~$\sFmla{\False}{!A}$: ``يجوز أن تكون~$!A$ كاذبة''،
وذلك في بعض \iftag{FOL}{!!{structure}}{!!{valuation}}.

وكل !!{signed formula} في الشجرة لا تخلو من أحد أمرين: إما أن تكون
\emph{افتراضًا}، ومكان الافتراضات أعلى الشجرة، وإما أن تكون مأخوذة من
!!a{signed formula} تعلوها بإحدى قواعد الاستدلال. ولكل !!{main operator}
قد يتصدّر ال!!{formula} السابقة قاعدتان؛ إحداهما لإشارة~$\True$،
والأخرى لإشارة~$\False$. فمن القواعد ما يشعّب الشجرة، ومنها ما يقتصر
على إلحاق !!{signed formula}s بالفرع نفسه. وليس من شرط القاعدة أن
تتناول ال!!{signed formula} السابقة لها مباشرة؛ بل يجوز لها، وكثيرًا
ما يجب، أن تتناول أي !!{signed formula} على مسار الفرع، من الجذر
إلى الموضع الذي تجرى فيه القاعدة.

والفرع \emph{مغلق} متى اجتمع فيه
$\sFmla{\True}{!A}$ و$\sFmla{\False}{!A}$؛ فإذا أغلقت الفروع كلّها
كان ال!!{tableau} مغلقًا. ووجه ذلك في التفسير الحدسي أن الفرع يصف
حالًا يُفترض إمكان اجتماع أجزائها، لكن~$\sFmla{\True}{!A}$
و$\sFmla{\False}{!A}$ لا يجتمعان. فإغلاق الفرع إبطال للإمكان الذي
كان يصفه. ومن ثم فإن ال!!{tableau} المغلق يبطل جميع الوجوه التي
قد يصدق فيها كل افتراض صورته~$\sFmla{\True}{!A}$، ويكذب فيها كل
افتراض صورته~$\sFmla{\False}{!A}$، معًا.

ونسمّي ال!!{tableau} المغلق \emph{لـ~$!A$} كل !!{tableau} مغلق
جذره~$\sFmla{\False}{!A}$. فإذا وجد !!{tableau} مغلق على هذا الوجه،
بطلت جميع إمكانات كذب~$!A$؛ فلا بد أن تصدق~$!A$ في كل
\iftag{FOL}{!!{structure}}{!!{valuation}}.

\end{document}
